\section{Evaluate the causal contract across the agent's components}
\label{exec18:framework-evaluation}

An agent framework becomes useful to the bank when its components can be
evaluated against a stable decision contract. Consider again an arrears
servicing offer. A plausible explanation, a completed sequence of tool calls
and a correct causal answer are three different outcomes. The system might
produce a fluent justification while conditioning on a variable recorded
after the offer. It might execute the requested estimator successfully while
answering the effect among treated accounts rather than the effect in the
eligible population. Conversely, a scientifically correct answer might be
that the declared observations do not identify the requested effect. A
framework comparison should preserve these distinctions when it chooses
components, evaluation data and metrics.

\subsection{Architecture is a map of responsibilities, not a certificate}

\citet{exec18_framework_agents_survey} organize autonomous-agent architecture
into profile, memory, planning and action, and distinguish architecture from
capability acquisition. This is a useful account of the survey authors'
taxonomy. It is secondary evidence for the individual systems and experiments
that the survey discusses. Its selected March 2025 revision, including a
benchmark table largely describing 2022--2023 work, cannot establish the
current frontier through October 2026. The taxonomy nevertheless helps
specify where a causal assistant's responsibilities reside.

A profile can declare a servicing analyst's remit; it cannot establish the
bank's authority to change an account's treatment. Memory can preserve a
previous graph and its supporting sources; retrieval of that graph does not
make its assumptions current for another product or calendar period.
Planning can propose an adjustment analysis, a sensitivity analysis or a
request for clarification. The plan needs a mathematical applicability
check before execution. An action component can call a causal engine, but
the result must retain the target, graph and observation definitions used
by that engine. These are proposed mappings from the taxonomy to banking
responsibilities, not findings that the surveyed agents already satisfy
them. The state invariant in Equation~\ref{exec14:l6-state} remains the
application's responsibility.

The survey separates subjective human evaluation from objective metrics,
protocols and benchmarks. Task success, similarity to human behavior and
efficiency answer different questions. Its discussion also recognizes the
cost and population bias of human evaluation and the inability of objective
techniques to measure every capability perfectly. In this proposal, a
business reviewer judges whether the question and explanation reflect the
intended decision. A mathematical reference checks whether a result is
correct under explicit assumptions. Neither judgment should silently
replace the other. An uncertain or disputed reference requires an
adjudication status; another model's agreement is not sufficient causal
ground truth. \citep{exec18_framework_agents_survey}

The same survey identifies interacting-module prompt fragility,
hallucination, knowledge-boundary problems in simulations and repeated-call
efficiency. These motivate concrete evaluation cases. A retrieved passage
can contain an instruction inconsistent with the declared target. An
apparently helpful memory can introduce information unavailable at the
decision time. A simulated analyst can exploit knowledge hidden from the
person being represented, making the simulation's task success irrelevant
to the real information constraint. A correct final answer can require many
failed intermediate calls. These are reasons to inspect complete traces and
information availability, rather than infer banking performance from a
general agent benchmark. The cases are proposed causal evaluations;
their frequencies have not been measured here.
\citep{exec18_framework_agents_survey}

\subsection{What a general evaluation library contributes}

\citet{exec18_framework_promptbench} present PromptBench as a modular library
for models, datasets, prompts, adversarial attacks and evaluation protocols.
Their four-stage pipeline initializes the data and model, specifies prompts,
defines input/output handling and metrics, and runs the evaluation. The
selected article reports twelve tasks and twenty-two datasets. Its appendices
describe extension hooks for datasets, models, methods and metrics. These
are relevant design mechanisms for a causal-contract evaluator: a dataset
could contain versioned business questions and admissible responses, while
a custom metric checks target and assumption preservation.

Those mechanisms do not supply the causal reference. For example, a case
asking for an adjustment set needs a declared graph and an appropriate
graphical criterion. A case asking for an effect interval needs a reference
identified set or another justified numerical target. A question whose
meaning changes with the timing of a balance variable needs a clarification
reference. Matching a stored phrase without checking these objects would
reward superficial agreement. A proposed adapter should therefore send a
typed request to the existing scientific boundary and score its typed
response, while retaining the prompt, reference version and execution
trace. PromptBench's extensibility makes that design plausible; no such
adapter is claimed to have been implemented or qualified by this source
inspection. \citep{exec18_framework_promptbench}

The paper's empirical conclusions also argue against a single preferred
prompt. It reports adversarial vulnerability, prompt-engineering methods
that do not outperform the baseline on every dataset, and remaining
weaknesses on selected dynamic tasks. These are attributed historical
results, not a local reproduction or a ranking of present models. Its own
limitations discuss incomplete scenario coverage, metric nuance and the
quality of datasets and prompts. A banking comparison should consequently
retain the unmodified baseline, prespecify perturbations and report
task-specific failures. It should not select a favorable wording after
examining held-out answers and then call that result general robustness.
The paper's historical comparison with other libraries is likewise an
author comparison at that time, not an independent purchasing table.
\citep{exec18_framework_promptbench}

The proposed reuse disposition is reference use of these evaluation
interfaces, followed by conditional integration if a selected implementation
fits the causal request and result schemas. A pinned code revision, package
and dependency rights, environment, input/output behavior and a matched
baseline would have to be inspected before executable adoption. The
article's CC-BY-4.0 notice concerns the paper; it is not evidence of the
software distribution's license. Porting the full library would require
new evidence that its maintenance cost is justified relative to a small
local evaluator. Rejecting every general evaluator merely because it does
not itself prove identification would also discard potentially useful
infrastructure. The appropriate boundary is to reuse orchestration and
measurement only while retaining application-owned scientific references.

\subsection{Completion, correctness and the denominator}

Evaluation must make unsuccessful execution visible without penalizing a
correct scientific refusal as if it were a crash. For a fixed case $i$ and
wording $j$, let $R_{ij}=1$ denote a completed, interpretable response and
$R_{ij}=0$ an unanswered execution. Among completed responses, let
$K_{ij}=1$ denote correctness under the frozen causal contract and
$K_{ij}=0$ an incorrect response. Set $R_{ij}K_{ij}=0$ for an unanswered
case solely for the following accounting identity. With $N$ planned
requests, define
\begin{equation}
 C=\frac{\sum_{ij}R_{ij}}{N},\qquad
 J=\frac{\sum_{ij}R_{ij}K_{ij}}{N},\qquad
 A=\frac{\sum_{ij}R_{ij}K_{ij}}{\sum_{ij}R_{ij}},
 \qquad J=CA .
 \label{exec18:framework-denominators}
\end{equation}
Here $C$ is completion, $J$ is correct responses per planned request and
$A$ is correctness conditional on completion. The last two expressions
involving $A$ apply only when at least one response is completed; $A$ is
undefined if none is completed. These are finite accounting quantities,
not estimates of calibrated causal uncertainty. A correct
nonidentification certificate, an appropriate ambiguity report or a
required clarification can count as completed and correct. The contract
must declare when each is appropriate; indiscriminate refusal cannot
receive credit on answerable cases. \citep{exec18_framework_reference}

An exact authored example makes the difference material. Stipulate twenty
servicing-question families, each with two wordings having identical
outcomes within that family for each system. There are forty planned
requests. The example specifies the answers directly; it contains no model
calls, framework trials, simulated accuracy or bank observations.
\begin{table}[htbp]
\centering
\small
\begin{tabular}{@{}lrrrrrr@{}}
\toprule
System & Correct & Wrong & Unanswered & $C$ & $J$ & $A$\\
\midrule
A & 22 & 2 & 16 & $3/5$ & $11/20$ & $11/12$\\
B & 30 & 6 & 4 & $9/10$ & $3/4$ & $5/6$\\
\bottomrule
\end{tabular}
\caption{Exact stipulated outcomes on forty planned requests in twenty
paired families. Conditional correctness favors A; correct responses per
planned request favor B. No sampling or model-performance claim is made.}
\label{exec18:framework-example}
\end{table}
Thus reporting $A$ alone favors the system that completes fewer cases,
whereas reporting $J$ alone conceals its lower error count. Neither scalar
settles the banking decision. The full counts distinguish reduced coverage
from increased wrong answers, and task strata can reveal whether failures
concentrate on nonidentification, timing or diagram construction.
\citep{exec18_framework_reference}

The paired design also limits uncertainty claims. Two wordings of one
contract are not two new scientific targets. In this example their outcomes
are stipulated to be identical, so treating all forty requests as
independent Bernoulli observations would invent replication. For a future
sampled evaluation, comparisons should retain pairing across systems and
the dependence among wordings of a contract. The sampling law over
contracts, shared model behavior and adjudicator uncertainty determine
whether a cluster analysis, a paired randomization argument or another
method is justified. This fixed example supplies no such sampling law and
therefore supplies no confidence interval.

\subsection{Why the cost of an error can reverse the choice}

Suppose, strictly for illustration, that a correct response has value
$v=10$ USD, a wrong response has loss $p$ USD, and unanswered requests have
zero value before their execution costs. Let $n_{+}$ and $n_{-}$ be the
correct and wrong counts, and let $c$ be total execution and review cost in
USD. Net value per planned request is then
\begin{equation}
 V(p)=\frac{v n_{+}-p n_{-}-c}{N}.
 \label{exec18:framework-value}
\end{equation}
Stipulate total costs $c_A=12$ USD and $c_B=24$ USD over the forty requests.
For $p=20$ USD, $V_A=21/5=4.20$ USD and
$V_B=39/10=3.90$ USD per planned request. For $p=5$ USD,
$V_A=99/20=4.95$ USD and $V_B=123/20=6.15$ USD.
The same stipulated answers thus favor different systems under different
error losses. Every amount is an assumption or exact arithmetic consequence;
none is an observed bank price, benefit or return on investment.
\citep{exec18_framework_reference}

A real servicing evaluation would also need to value justified refusal,
clarification, delay and downstream human work, and distinguish errors
with different monetary consequences. A spurious recommendation to change
a customer's treatment is not interchangeable with an incorrect citation.
Those values must be elicited for a declared decision and information set.
Framework token or request counters contribute to measured execution cost;
they cannot establish these benefits or losses. A low-cost system can be
unattractive if its harmful errors exceed the acceptable rate, even when
the average in Equation~\ref{exec18:framework-value} is positive.

The practical comparison is therefore a paired, versioned set of causal
contracts with explicit completion states, scientific references, trace
inspection and a separately stated economic rule. General agent taxonomies
help locate responsibilities, and general evaluation libraries can help
run that comparison. They leave causal truth, reference construction,
rights qualification and the justified sampling law to be established.
The proposed architecture remains conditional on that work; these two
originals do not establish exhaustive framework coverage or a qualified
banking agent.
